% =====================================================================
%  Chapter 37 — The Law
% =====================================================================

\chapter{The Law}
\label{ch:law}

The owner asked for every legal advantage. This chapter sets out the legal
envelope the system operates in, and it is short, because the envelope is
wide and the system stays well inside it.

\section{What the reader is doing}

The reader buys and sells exchange-traded funds in their own account, for
their own money, without telling anyone what to do and without being paid
by anyone to do it.

Under the Investment Advisers Act of 1940, a person who advises others
about securities for compensation must register as an investment adviser.
Section 9(a)(2) is the operative prohibition on doing so without
registration, and Section 9(a)(3) prohibits fraudulent, deceptive or
manipulative acts by an adviser. A person managing their own money is
outside both: the Act's definition of an investment adviser excludes
investment advisers who, as a substantial part of their regular business,
provide advice with respect to securities only to investment accounts
over which they exercise investment discretion, and only to accounts they
own, in each case without compensation.

The practical rule is short. Managing your own money is not advice to
anyone else and requires no registration. Running the same system on
behalf of a family member, a colleague or a paying client is a different
activity entirely, and it is the point at which this chapter stops being
relevant and a lawyer starts being relevant.

\begin{takeaway}[The one-line rule]
Trading your own account is not advisory activity and needs no
registration. Doing the same thing for someone else, for compensation or
with discretion over their money, is. This book assumes the first case and
says nothing useful about the second.
\end{takeaway}

\section{What the system does not do}

Three things a system of this kind could conceivably do that would
create legal exposure, and that this one does not.

\textbf{It does not manipulate.} Section 10(b) of the Securities Exchange
Act and Rule 10b-5 prohibit fraud in connection with the purchase or
sale of securities, and Rule 10a-1 and 10a-2 prohibit manipulative
practices. The system buys a fund because its model ranks that fund
favourably relative to the other funds available, and it sells a fund
because its model or the governor says to. It does not create a
synthetic quote, does not wash a trade between accounts, does not buy in
order to move a price, and does not have the market impact to move a
price in any case --- Chapter 11 showed that the reader's position is a
rounding error against the fund's daily volume.

\textbf{It does not disclose anything selectively.} Regulation FD, at
17~CFR 243.100 through 243.103, requires a company that is an issuer or
an insider to disclose material non-public information to securities
holders and to the market. A private investor who is not an insider and
who does not receive issuer information is not subject to it. The
provision matters to this book only as a warning: the moment a reader
begins to receive information from an issuer in connection with a
position, the analysis of this chapter stops applying.

\textbf{It does not borrow.} No part of the system uses margin, and
Chapter 38's recommendations are all unlevered. This matters because the
margin rules have changed recently and because borrowing turns a
drawdown into a forced sale, which is the one thing the entire architecture
of this book exists to prevent.

\section{The margin rules, and a flag}

The Financial Industry Regulatory Authority issued Regulatory Notice
26-10 on 20 April 2026, effective 4 June 2026, which replaced the
day-trading margin requirements --- including the familiar \$25,000
minimum equity for a pattern day trader --- with new intraday margin
standards under Rule 4210(d)(2).

The book's environment presents one complication that must be flagged
rather than smoothed over: the Securities and Exchange Commission's
investor education page on day trading still describes the \$25,000
pattern-day-trader requirement. The two are not consistent, and a reader
who relies on the investor education page rather than the regulatory notice
would be relying on a statement that the notice has superseded.

The system in this book is unaffected either way, because it does not day
trade, does not use margin, and trades monthly. The flag is included
because the book is about taking every advantage, and an advantage that
consists of acting on stale regulatory guidance is not an advantage.

\section{The tax rules, as compliance rather than as arithmetic}

The tax layer of Chapter 29 is not only an arithmetic model. It is a
statement of what the reader must actually do.

The wash-sale rule is a reporting and substantiation requirement: the
basis of the replacement shares is adjusted, the holding period carries
over, and the disallowed loss is tracked. A reader running this system
must keep the lot records that make that possible. The book's engine
keeps them, and the data they need --- the acquisition date, the
quantity, the basis, the disposal date --- is exactly the data a broker's
tax lot report provides.

The capital-loss carryforward and the \$3,000 annual offset against
ordinary income are also substantiation items. The engine computes them,
but the reader must report them.

\section{Exchange-traded funds, and the investor bulletin}

The Securities and Exchange Commission's investor bulletin on exchange-
traded funds makes two points that bear directly on this book's costs.
The first is that the bid--ask spread is a real cost to the investor that
does not appear in the expense ratio, and that most funds publish their
median spread on their own website. The second is that an ETF's price can
diverge from the value of the assets it holds.

Neither point creates a legal exposure, and both are handled. The spread
is charged in the book's cost model. The divergence is a consequence of
the creation and redemption mechanism: an authorised participant
exchanges a basket of securities for ETF shares and can arbitrage the
difference, which is what keeps the price close to the value.

\begin{takeaway}[Why the ETF structure is the safe wrapper]
The creation and redemption mechanism is what keeps an ETF's price close
to the value of its holdings, and it is the reason a single investor
cannot move an ETF's price against itself. A single investor buying a
fraction of a per cent of a fund's daily volume is not manipulating
anything; a single investor cannot.
\end{takeaway}

\section{The limits of this chapter}

Three exclusions, stated rather than hidden.

This is not legal advice, and a reader with specific circumstances should
obtain it. This chapter is a description of the envelope as it stood at
the date on the copyright page, not a legal opinion.

It does not cover state law, which varies on the treatment of some
securities and on the definition of advice.

And it does not cover the reader's own employment. A reader who is
subject to a restricted list, a personal account dealing policy, or a
conflict-of-interest rule at their employer has obligations that this
chapter does not know about, and those obligations are strict.

\begin{keytakeaways}
\item Trading your own account requires no registration; advising others
for compensation does, and this book assumes the first case.
\item The system does not manipulate, does not make selective disclosures,
and does not borrow --- the last of which is what keeps a drawdown from
becoming a forced sale.
\item The margin rules changed in June 2026 and the investor education
material has not caught up; the book flags the conflict rather than
reconciling it.
\item The tax layer is a compliance record as well as an arithmetic model,
and the wash-sale bookkeeping is a substantive obligation.
\end{keytakeaways}
